\subsection{COMPLEMENTS OF UNIONS AND INTERSECTIONS}

PROPOSITION 5. For every family \((\mathbf{X}_t)_{t\in I}\) of subsets of a set E, we have

\[
\mathbb {G} _ {\mathrm{E}} \left(\bigcup_ {i \in \mathrm{I}} \mathrm{X} _ {i}\right) = \bigcap_ {i \in \mathrm{I}} \left(\mathbb {G} _ {\mathrm{E}} \mathrm{X} _ {i}\right), \quad \mathbb {G} _ {\mathrm{E}} \left(\bigcap_ {i \in \mathrm{I}} \mathrm{X} _ {i}\right) = \bigcup_ {i \in \mathrm{I}} \left(\mathbb {G} _ {\mathrm{E}} \mathrm{X} _ {i}\right).
\]

Let \(x \in \complement_{\mathbf{E}} \left( \bigcup_{i \in \mathbf{I}} \mathbf{X}_i \right)\) . Then \(x \in E\) and, for every \(i \in I\) , \(x \notin \mathbf{X}_i\) , so that \(x \in \complement_{\mathbf{E}} \mathbf{X}_i\) ; consequently

\[
x \in \bigcap_ {i \in I} \left(\complement_ {E} X _ {i}\right).
\]

Conversely, let \(x\) be an element of \(\bigcap_{i \in I} (\complement_E X_i)\) ; by the definition of intersection we have \(x \in E\) . Furthermore, if we had \(x \in \bigcup_{i \in I} X_i\) , there would exist an index \(x \in I\) such that \(x \in X_x\) , which contradicts the hypothesis \(x \in \bigcap_{i \in I} (\complement_E X_i)\) ; hence

\[
x \in \mathbb {C} _ {\mathbf {E}} \left(\bigcup_ {i \in \mathbf {I}} \mathrm{X} _ {i}\right).
\]

This proves the first formula; the second one is an immediate consequence, in view of the relation \(\mathbf{G}_{\mathbf{E}}(\mathbf{G}_{\mathbf{E}}\mathbf{X})=\mathbf{X}\) for every subset X of E.

